\section{The joint gap and norm certificate}
\label{joint:section}

\begin{theorem}[Sharp joint envelope]\label{joint:main}
For $g,P>0$, put
\begin{equation}\label{joint:definition}
 h=\min\{g,P^2/m\},\qquad
 \mathcal J(g,P)=
 \max_{h^2/P^2\le\rho\le h/m}V(h,\rho).
\end{equation}
Set $\mathcal J(g,0)=\mathcal J(0,P)=0$. Under
\eqref{boundarygap:assumptions},
\begin{equation}\label{joint:bound}
 \delta(x)\le\mathcal J(G(x),P(x))
       \le\min\{\mathcal B(G(x)),\mathcal E(P(x))\}.
\end{equation}
The function $\mathcal J$ is finite, continuous on the closed quadrant,
and nondecreasing in each argument. If $g\ge P^2/m$, the interval
collapses and $\mathcal J(g,P)=\mathcal E(P)$.

For every $g,P>0$, $\mathcal J(g,P)$ is the supremum of the source
error with either the two upper bounds $G(x)\le g$, $P(x)\le P$, or
the two exact observations $G(x)=g$, $P(x)=P$. Smooth interval examples
suffice for upper bounds; exact observations need at most two dimensions.
Their compact convex domains have nonempty interior in their ambient space
and may vary.
\end{theorem}

\begin{proof}\label{joint:proof}
Consider upper bounds $g,P>0$ and a nonoptimal source of radius $r$.
Its terminal radial slope is at most $\min\{P,g/r\}$ and
$r\le\min\{P/m,\sqrt{g/m}\}$. If $g\ge P^2/m$, the norm theorem
gives the claim. Otherwise let $r_0=g/P$ and $R=\sqrt{g/m}$.
For $r\le r_0$, Lemma~\ref{compact:radial} gives
\[
 \delta(x)\le A(P,r)\le A(P,r_0)=V(g,r_0^2).
\]
For $r\ge r_0$, the terminal cap $g/r$ gives
$\delta(x)\le V(g,r^2)$. Combining the ranges gives
\eqref{joint:definition}. Each fixed-radius interval construction
meets both upper bounds, so its loss is bounded by both separate
envelopes. In the norm branch the stationary interval has gap
$P^2/m\le g$. These constructions also prove sharpness for upper bounds.

To keep both observations exact when $g<P^2/m$, choose a maximizing
$\rho$, set $r=\sqrt\rho$, $p=g/r\le P$, $q=p-mr$, and use smooth
near-extremizers $u_i$ with initial slope $q$, terminal slope $p$ and
$u_i(r)\to V(g,\rho)$. For any $a>0$, set
\[
 c=\sqrt{P^2-p^2},\qquad Q=[0,r]\times[0,a],\qquad
 f_i(s,y)=u_i(s)+cy+my^2/2.
\]
The unique minimum is $(0,0)$, and
\[
 \langle\nabla f_i(s,y)-\nabla f_i(0,0),(s,y)\rangle
 =(u_i'(s)-q)s+my^2\ge m(s^2+y^2).
\]
The Hessian is diagonal, with norm
$\max\{|u_i''(s)|,m\}\le H_0+H_1u_i(s)\le H_0+H_1f_i(s,y)$.
At $(r,0)$ the gradient is $(p,c)$, of norm $P$. Its components are
nonnegative, so its full gap on $Q$ is exactly $rp=g$.
The source losses converge to $\mathcal J(g,P)$.

If $g\ge P^2/m$, let $r=P/m$ and $a=g/P-r\ge0$.
Use stationary smooth near-extremizers on $[0,r]$ and extend them to
$[-a,0]$ by the same positive-curvature solution with which they agree
near zero:
\[
 u_+(s)=
 \begin{cases}
 (H_0/H_1)\{\cosh(\sqrt{H_1}s)-1\},&H_1>0,\\
 H_0s^2/2,&H_1=0.
 \end{cases}
\]
This even extension is smooth, has its unique minimum at zero, and
satisfies $u_+'(s)s\ge H_0s^2\ge ms^2$ and the curvature bound with
equality. At $r$, the norm is $P$, the gap is $(r+a)P=g$, and the
loss tends to $\mathcal E(P)$. The quadratic case $H_0=m,H_1=0$ is
already smooth.

The upper-bound characterization proves monotonicity. For positive
$(g,P)$, parameterize the compact interval in
\eqref{joint:definition} by $z\in[0,1]$; continuity of $V$ gives
continuity of the maximum, including its collapse.
At the axes use $0\le\mathcal J(g,P)\le\min\{\mathcal B(g),\mathcal E(P)\}$.
If either upper bound is zero, the source is optimal.
Exact data with $g=0$ and any $P\ge0$ occur at the left endpoint of
$f(s)=Ps+ms^2/2$. Exact data with $P=0<g$ are infeasible; the stated
zero extension there concerns upper bounds only.
\end{proof}

\begin{corollary}[Strict improvement]\label{joint:strict}
For $H_1\ge0$ and $0<g<P^2/m$,
\[
 \mathcal J(g,P)\le\mathcal E(P)
   -\frac{(P^2-mg)^2}{2mP^2}<\mathcal E(P).
\]
If also $H_1>0$ and $\mathcal S(g)\le0$, then
$\mathcal J(g,P)<\min\{\mathcal B(g),\mathcal E(P)\}$.
\end{corollary}
\begin{proof}
For $\rho\in[g^2/P^2,g/m]$, put $p=g/\sqrt\rho$ and
$q=p-m\sqrt\rho$. Apply \eqref{bndnorm:radial-deficit} to smooth
near-extremizers, and use \eqref{compact:norm-growth}:
\[
 \mathcal E(P)-V(g,\rho)
 \ge\frac{P^2-p^2+q^2}{2m}
 =\frac{P^2-2mg+m^2\rho}{2m}
 \ge\frac{(P^2-mg)^2}{2mP^2}.
\]
Maximization proves the first claim. Under the second set of assumptions,
$V(g,\rho)<g$ on the compact joint interval, whose left endpoint is
positive. Thus its maximum is below $g=\mathcal B(g)$.
\end{proof}
The displayed norm deficit is attained when $H_0=m$ and $H_1=0$.

\begin{corollary}[One root or an algebraic bound]
\label{joint:ordinary}\label{joint:regime}
Let $H_1>0$, $g,P>0$, and $h=\min\{g,P^2/m\}$.
If $\mathcal S(h)\le0$, then
\begin{equation}\label{joint:ordinary-formula}
 \mathcal J(g,P)=V(h,h^2/P^2)
 \le h+\frac{h^2}{P^2}\mathcal S(h)\le h.
\end{equation}
\end{corollary}
\begin{proof}
Apply the radius monotonicity and tangent bound from
Corollary~\ref{gapregime:radius-monotone} at the left endpoint of the
joint interval. This includes a collapsed interval.
\end{proof}
The exact envelope needs one switch root. The algebraic upper bound
needs only the square root in $\mathcal S(h)$ and arithmetic, but
does not promise a prescribed relative approximation of $\mathcal J$.
For $m=H_0=H_1=1$, $g=1$, $P=2$, it gives
\[
 \mathcal S(1)=4-\frac{7\sqrt{14}}6<0,\qquad
 \mathcal J(1,2)\le2-\frac{7\sqrt{14}}{24}<1.
\]

\begin{corollary}[Constant curvature]\label{joint:constant}
If $H_1=0$, then for $g,P>0$,
\begin{equation}\label{joint:constant-formula}
 \mathcal J(g,P)=
 \begin{cases}
 h+\beta h/m,&\beta\ge0,\\
 h+\beta h^2/P^2,&\beta<0,
 \end{cases}
 \qquad h=\min\{g,P^2/m\}.
\end{equation}
If $m\le H_0<(1+\sqrt2)m$ and $0<g<P^2/m$, it is strictly below
both separate envelopes.
\end{corollary}
\begin{proof}
Maximize the affine function $V(h,\rho)=h+\beta\rho$ at the appropriate
interval endpoint. For $\beta<0$ and $h=g$ the result is below
$g=\mathcal B(g)$; Corollary~\ref{joint:strict} gives the norm comparison.
\end{proof}
For example, $m=H_0=P=1$, $H_1=0$, $g=1/2$ give
$\mathcal J(1/2,1)=3/8$, while $\mathcal B(1/2)=\mathcal E(1)=1/2$.
The function $f(s)=s/2+s^2/2$ on $[0,1/2]$ attains the joint value.

For finite evaluation, apply Proposition~\ref{joint:budget} with gap
argument $h$, interval $[h^2/P^2,h/m]$ and $\ell=h/2$.
The norm branch uses a single node; the ordinary-gap regime uses its
left node; otherwise use the interval grid.
The lower bound $\ell$ concerns the extremal envelope, not the actual error.

Certified upper observations $G(x)\le g^+$, $P(x)\le P^+$ give the
valid lower bound $f(x)-\mathcal J(g^+,P^+)$ on the optimum.
A bound on feasible-direction score errors alone does not supply $P^+$.
When both observations are already available, the joint formula uses no
additional objective, gradient or linear minimization oracle query.
Its scalar evaluation remains a separate cost.
Since $\mathcal J(g^+,P^+)\le\mathcal B(g^+)$, accumulating these lower
bounds preserves the gap-only work estimates of Theorem~\ref{nloracle:main};
no inverse of a two-variable function is needed.
